\ProvidesFile{FunctionList.tex}[Списки функций]


\subsection{Списки функций}
\label{section::FunctionList}

\subsubsection{Что мы хотим}

Здесь я хочу обсудить, как работать работать со списком функций во время компиляции.
Звучит супер просто, учитывая, что мы уже знаем, как хранить произвольные элементы в списке.
Действительно
\begin{cppcode}
namespace AL {
template<auto...>
struct List {};
}

void f(int) {}

struct A {
  void f() const {}
};

int main() {
  using List = AL::List<&f, &A::f>;
  AL::At<List, 0>(1); // calls f(1)
  A a;
  (a.*AL::At<List, 1>)(); // calls a.f()
}
\end{cppcode}
Однако, из-за того, что в языке есть понятие перегрузки функций и операторов, то одним и тем же именем могут называться разные методы.
Например в следующих ситуациях у нас будет ошибка компиляции.
\begin{cppcode}
namespace AL {
template<auto...>
struct List {};
}

void f(int) {}

void f(double) {}

struct A {
  void f() const {}
  void f() {}
};

int main() {
  using List = AL::List<&f>; // Compilation error ambiguity
  using List = AL::List<&A::f>; // Compilation error ambiguity
}
\end{cppcode}
Проблема в том, что компилятор не знает, какую из перегруженных функций вы имеете в виду в этом месте.
Для решения этой проблемы нужно указать явно сигнатуру функции.
Это делается с помощью \verb"static_cast" следующим образом
\begin{cppcode}
using List = AL::List<static_cast<void(*)(int)>(&f)>;
using List = AL::List<static_cast<void (A::*) const>(&A::f)>; 
\end{cppcode}
Как можно заметить в первом случае синтаксис содержит неудобную для использования \verb"(*)", которая не несет информации для программиста, но необходима.
А во втором случае мы дважды используем имя класса \verb"A": один раз при указании сигнатуры, а второй раз при обращении к адресу метода.
Наша задача упросить процесс обращения к перегруженным методам.
Мы бы хотели иметь следующий синтаксис:
\begin{cppcode}
using List = AL::List<sig<void(int)>(&f)>; // selects f(int)
using List = AL::List<sig<void () const>(&A::f)>;  // selects A::f() const
\end{cppcode}

\subsubsection{Имплементация \texttt{sig}}

Прежде всего начнем с простой шаблонной имплементации:
\begin{cppcode}
namespace detail {
template<class signature>
struct Sig {
  static_assert(Logic::False<signature>,
                "The argument of sig MUST be a valid function signature!");
};

template<class R, class... Args>
struct Sig<R(Args...) const> {
  using consnt_signature = R(Args...) const;

  template<class T>
  constexpr consnt_signature T::*
  operator()(consnt_signature T::*f) const noexcept {
    return static_cast<consnt_signature T::*>(f);
  }
};

template<class R, class... Args>
struct Sig<R(Args...)> {
  using signature = R(Args...);

  template<class T>
  constexpr signature T::*operator()(signature T::*f) const noexcept {
    return static_cast<signature T::*>(f);
  }
};
} // namespace detail

template<class signature>
detail::Sig<signature> sig = {};
\end{cppcode}
Давайте пройдемся по имплементации.
\begin{enumerate}
\item Прежде всего обратите, что \verb"sig" является не шаблоном класса, а шаблонной переменной.
А круглые скобки -- это перегруженный оператор.

\item Как работает оператор \verb"operator()".
Если пользователь ввел что-то вида
\begin{cppcode}[numbers = none]
sig<R(Args...) const>(&A::f);
\end{cppcode}
В этом случае
\begin{cppcode}[numbers = none]
signature = R(Args...) const
\end{cppcode}
А получение из этой сигнатуры указателя на объект класса \verb"T" делается так:
\begin{cppcode}[numbers = none]
signature T* = R (T::*)(Args...) const
\end{cppcode}
Это просто возможность языка, о которой надо знать.

\item После чего оператор \verb"operator()" принимает только функции с сигнатурой
\begin{cppcode}[numbers = none]
R (T::*)(Args...) const
\end{cppcode}
для любых классов \verb"T".
Это позволяет избавиться от необходимости указывать явно тип класса.

\item Аналогичная имплементация для не константного метода.
Так же заметим, что типы
\begin{cppcode}[numbers = none]
R (T::*)(Args...) const
R (T::*)(Args...)
\end{cppcode}
в C++ являются разными и их имплементации нельзя объединить в одну.
Приходится для каждой из них писать отдельно.

\item Причина использования оператора \verb"operator()" следующая.
Мы можем задать операцию на элементах по сути двумя способами:
\begin{enumerate}
\item Шаблоном
\begin{cppcode}[numbers = none]
template<auto arg>
struct FImpl {
  static constepxr auto value = /*implementation*/;
};
\end{cppcode}

\item С помощью \verb"constexpr" функции
\begin{cppcode}[numbers = none]
template<class T>
constexpr auto function(T arg) { 
  return /*implementation*/ 
}
\end{cppcode}
\end{enumerate}
Важная особенность второго подхода в том, что мы можем иметь много разных функций с одним именем, но разной сигнатурой.
Кроме того, для аргументов функции недостающие шаблонные параметры определяются автоматически.
Это позволяет не указывать явно класс \verb"T", которому принадлежит метод \verb"f".
\end{enumerate}
Теперь мы действительно можем писать код вида
\begin{cppcode}
namespace AL {
template<auto...>
struct List {};
}

void f(int) {}

void f(double) {}

struct A {
  void f() const {}
  void f() {}
};

int main() {
  using List = AL::List<sig<void(int)>(&f), sig<void(double)>(&f), 
                        sig<void() const>(&A::f), sig<void()>(&A::f)>;
}
// List equals AL::List<f(int), f(double), A::f() const, A::f()>
\end{cppcode}

\subsubsection{Смягчение требований по \texttt{const}}

Давайте рассмотрим следующий набор классов:
\begin{cppcode}
struct A {
  void f() const {}
  void f() {}
};

struct B {
  void f() {}
};

struct C {
  void f() const {}
};
\end{cppcode}
Какую из функций можно позвать по сигнатуре ниже?
\begin{cppcode}[numbers = none]
void() const
\end{cppcode}
Только константные методы \verb"A::f() const" и \verb"C::f() const".
Позвать \verb"A::f()" и \verb"B::f()" не получится, так как мы не можем менять константный объект.
С другой стороны, какой из методов мы можем вызвать по сигнатуре такой?
\begin{cppcode}[numbers = none]
void()
\end{cppcode}
На самом деле тут можно вызвать любой из $4$ методов.
Это согласуется с тем, что константный метод можно вызывать как на константном, так и на не константном объекте.
В этом случае мы бы хотели, чтобы сигнатура
\begin{cppcode}[numbers = none]
void() const
\end{cppcode}
означала выбираем только константный метода, а сигнатура
\begin{cppcode}[numbers = none]
void()
\end{cppcode}
означает, выбери не константный метод, если он есть, а если нет, то выбери константный метод.
То есть мы хотим такое поведение:
\begin{cppcode}
using List1 = AL::List<sig<void()>(&A::f), sig<void()>(&B::f), sig<void()>(&C::f)>;
// List1 equals AL::List<A::f(), B::f(), C::f() const>

using List2 =
    AL::List<sig<void() const>(&A::f), sig<void() const>(&B::f), sig<void() const>(&C::f)>;
// Compilation error because there is no const B::f method
\end{cppcode}
В этом случае имплементация для константной сигнатуры не изменится.
А вот для не константной сигнатуры надо будет написать так:
\begin{cppcode}

template<class R, class... Args>
struct Sig<R(Args...)> {
  using signature = R(Args...);
  using const_signature = R(Args...) const;

  template<class T>
  constexpr const_signature T::*operator()(const_signature T::*f) const volatile noexcept {
    return static_cast<const_signature T::*>(f);
  }

  template<class T>
  constexpr signature T::*operator()(signature T::*f) const noexcept {
    return static_cast<signature T::*>(f);
  }
};
\end{cppcode}
При этом вспомним, что
\begin{cppcode}
template<class signature>
detail::Sig<signature> sig = {};
\end{cppcode}
Давайте разберем как работает следующий код
\begin{cppcode}
constexpr auto f1 = sig<void()>(&A::f);
constexpr auto f2 = sig<void()>(&B::f);
constexpr auto f3 = sig<void()>(&C::f);
\end{cppcode}
В каждом из этих случаев вызывается оператор \verb"operator()", у которого два аргумента: неявный	 \verb"this" и явный -- адрес метода.
То есть для компилятора видит вызов вроде такого:
Давайте разберем как работает следующий код%
\footnote{Я специально не пишу название пространства имен \verb"detail", чтобы сделать код читаемым.
По хорошему конечно же надо писать \verb"detail::Sig<void()>" и так далее.}
\begin{cppcode}
constexpr auto f1 = Sig<void()>::operator()(&sig<void()>, &A::f);
constexpr auto f2 = Sig<void()>::operator()(&sig<void()>, &B::f);
constexpr auto f3 = Sig<void()>::operator()(&sig<void()>, &C::f);
\end{cppcode}

\paragraph{класс \texttt{A}}

Давайте явно запишем какие операторы получаются в классе \verb"Sig<void()>" после определения шаблонного параметра \verb"T = A".
Компилятор будет видеть что-то вроде такого:%
\footnote{Написать \texttt{void (A::)()} конечно не получится, но это идейная запись, для сокращения кода, чтобы он стал понятнее.}
\begin{cppcode}[numbers = none]
using const_signature = void (A::)() const;
const_signature* operator()(const volatile Sig<void()>*, const_signature*)
// and
using signature = void (A::)();
signature operator()(const Sig<void()>*, signature*)
\end{cppcode}
И пару \verb"sig<void()>" и \verb"&A::f" можно подставить и в первую и во вторую версию оператора
Однако, по определению \verb"sig<void()>" является \verb"const", но не \verb"volatile".
А значит для подстановки во вторую версию нужно использовать меньше кастов.
Во вторую версию можно сразу подставлять, а для первой версии нужно сделать каст к \verb"volatile".
Потому компилятор выбирает вторую версию.

\paragraph{класс \texttt{B}}

В в случае с классом \verb"B" все проще.
Компилятор видит единственную функцию \verb"f" с сигнатурой
\begin{cppcode}[numbers = none]
using signature = void (B::)();
\end{cppcode}
Которая подходит только под оператор
\begin{cppcode}[numbers = none]
signature* operator()(const Sig<void()>*, signature*)
\end{cppcode}
А значит будет выбран не константный метод \verb"B::f()".

\paragraph{класс \texttt{C}}

И наконец последний третий вызов для класса \verb"C".
Тут компилятор видит наоборот только константный метод с сигнатурой
\begin{cppcode}[numbers = none]
using const_signature = void (C::)() const;
\end{cppcode}
И он подходит только под первый оператор:
\begin{cppcode}[numbers = none]
const_signature* operator()(const volatile Sig<void()>*, const_signature*)
\end{cppcode}
Ибо второй нарушает \verb"const"-корректность.
И так как неявный каст к \verb"volatile" разрешен в языке, то выбирается первая версия оператора.
А потому возвращаемое значение будет единственным константным методом \verb"C::f() const".

\subsubsection{Глобальные функции как методы}

При работе с объектами часто удобно трактовать глобальные функции, как методы класса, считая, что они принимают указатель (или ссылку) на объект в своем первом аргументе:
\begin{cppcode}
struct A {};

void f(const A* a, int x) { /*implementation*/ }
void f(A* a, int x) { /*implementation*/ }
void f(const A& a, int x) { /*implementation*/ }
void f(A& a, int x) { /*implementation*/ }
\end{cppcode}
Такая трактовка нам может понадобиться для имплементации стирающих типов.
А именно, если в классе не хватает метода и вы не можете изменить класс объекта, то добавление глобальной функции с особым первым аргументом решит эту проблему.
Чтобы было понятно, что имеется в виду, я рекомендую посмотреть в раздел~\ref{}, а сейчас я лишь хочу обсудить как сделать выбор перегруженных методов со смягчением требования по \verb"const" для первого аргумента (если он ссылка или указатель).
Для этого надо сделать дополнительные имплементации для следующих сигнатур:
\begin{cppcode}[numbers = none]
R(const Arg1*, Args...)
R(Arg1*, Args...)
R(const Arg1&, Args...)
R(Arg1&, Args...)
R(Arg1&&, Args...)
\end{cppcode}
Это покроет все возможные случаи.
Так же обратите внимание, что эти случае пересекаются с сигнатурами
\begin{cppcode}[numbers = none]
R(Args...)
\end{cppcode}
в случае хотя бы одного аргумента.
То есть вы можете написать
\begin{cppcode}[numbers = none]
void f(const A*) {}
sig<void(const A*)>::(&f);
// or
struct A {
  void f(const A*) {}
};
sig<void(const A*)>(&A::f);
\end{cppcode}
В обоих случаях должна корректно выбираться нужная функция без ошибки компиляции.
А это значит, что в сигнатурах, которые мы разбираем сейчас надо рассматривать и глобальные функции и методы классов.
Я опишу как имплементировать два случая
\begin{cppcode}[numbers = none]
R(const Arg1*, Args...)
R(Arg1*, Args...)
\end{cppcode}
Остальные делаются аналогично.
Надо лишь сказать, что для случая \verb"Arg1&&" не должно быть константной версии, так как нет смысла делать rvalue константой.
Начнем с константного указателя:
\begin{cppcode}

template<class R, class Arg1, class... Args>
struct Sig<R(const Arg1*, Args...)> {
  using const_signature = R(const Arg1*, Args...) const;
  using free_const_signature = R(const Arg1*, Args...);

  template<class T>
  constexpr const_signature T::*
  operator()(const_signature T::*f) const noexcept {
    return static_cast<const_signature T::*>(f);
  }

  constexpr free_const_signature* operator()(free_const_signature* f) const volatile noexcept {
    return static_cast<free_const_signature*>(f);
  }
};
\end{cppcode}
Здесь все просто: шаблонный оператор для метода класса и не шаблонный оператор для функции.
Для не константного аргумента будет
\begin{cppcode}
template<class R, class Arg1, class... Args>
struct Sig<R(Arg1*, Args...)> {
  using signature = R(Arg1*, Args...);
  using const_signature = R(Arg1*, Args...) const;

  using free_signature = R(Arg1*, Args...);
  using free_const_signature = R(const Arg1*, Args...);

  template<class T>
  constexpr const_signature T::*operator()(const_signature T::*f) const volatile noexcept {
    return static_cast<const_signature T::*>(f);
  }

  template<class T>
  constexpr signature T::*operator()(signature T::*f) const noexcept {
    return static_cast<signature T::*>(f);
  }

  constexpr free_const_signature* operator()(free_const_signature* f) const volatile noexcept {
    return static_cast<free_const_signature*>(f);
  }

  constexpr free_signature* operator()(free_signature* f) const noexcept {
    return static_cast<free_signature*>(f);
  }
};
\end{cppcode}
Здесь так же приоритет задается с помощью \verb"volatile".
Обратите внимание, что в случае метода класса, мы не ослабляем \verb"const" требования к первому аргументу, что видно из введенных сигнатур:
\begin{cppcode}
using signature = R(Arg1*, Args...);
using const_signature = R(Arg1*, Args...) const;
\end{cppcode}
Потому что в этом случае пользователь нам указал в точности нужный тип первого аргумента и константность может быть только в неявном аргументе, который определяется шаблонным параметром \verb"T".
